% ECONOMICS_PROBLEM: {"id": "I13-450", "title": "Insurance-network design", "jel_code": "I13", "source_id": "OP-0055"}
% !TeX program = xelatex
\documentclass[11pt,a4paper]{article}
\usepackage[margin=23mm]{geometry}
\usepackage{fontspec}
\setmainfont{Latin Modern Roman}
\usepackage{amsmath,amssymb,mathtools}
\usepackage{xcolor,enumitem,booktabs,longtable,array,xurl,hyperref}
\definecolor{muted}{HTML}{53636C}
\hypersetup{colorlinks=true,urlcolor=blue,linkcolor=blue}
\setlength{\parindent}{0pt}
\setlength{\parskip}{5pt}
\newcommand{\R}{\mathbb R}
\newcommand{\E}{\mathbb E}
\newcommand{\N}{\mathbb N}
\newcommand{\ind}{\mathbf 1}
\newcommand{\Var}{\operatorname{Var}}
\newcommand{\Cov}{\operatorname{Cov}}
\newcommand{\supp}{\operatorname{supp}}
\DeclareMathOperator*{\argmin}{arg\,min}
\DeclareMathOperator*{\argmax}{arg\,max}
\newcommand{\entrymeta}[3]{{\sffamily\small\color{muted}\textbf{\texttt{#1}}\quad|\quad #2\par #3\par}}
\newcommand{\Problemref}[1]{\texttt{#1}}
\newcommand{\JELref}[2]{\texttt{#2} (JEL \texttt{#1})}
\begin{document}
\section*{Insurance-network design}
\textbf{Problem ID:} \texttt{I13-450}\quad \textbf{JEL:} \texttt{I13}\par
% BEGIN ECONOMICS STATEMENT
\label{problem:OP-0055}
\label{atlas:prob:H5}
\noindent\textbf{Original atlas identifier:} H5; entry 55. \textbf{Field:} Health economics. \textbf{Original tractability label:} Medium.

\noindent\textbf{Classification:} Parameterized theoretical/empirical research agenda; not a certified unsolved theorem.

\subsection*{Definitions and assumptions}
Let insurers $\ell=1,\ldots,L$ choose provider sets $N_\ell\subseteq\mathcal H$, premiums, and contracts in a market with provider set $\mathcal H$. Patients choose plans and providers subject to health needs, geography, preferences, and information. A model $m$ specifies negotiated prices, provider capacities, adverse selection, and network/plan equilibrium selection. Breadth is the network itself, not merely its number of hospitals: access depends on distance, specialty, capacity, and continuity. Define $EV_i(\mathbf N;m)$ as monetary patient welfare including premiums and health/access effects, excluding profit changes separately assigned to owners, and $\Pi_a(\mathbf N;m)$ as insurer/provider profits net of real costs. Fix welfare weights and a public budget/financing rule; legally enforceable adequacy requirements define the feasible regulation set $\mathcal P$.

\subsection*{Formal research question}
\emph{Editorial formulation: the mathematical target below is not a theorem quoted from the references.}

For regulation $p\in\mathcal P$, with resulting selected network equilibrium $\mathbf N(p;m)$, identify
\[
 W_m(p)=\sum_iw_iEV_i(\mathbf N(p;m);m)+\sum_aw_a\Pi_a(\mathbf N(p;m);m)-\lambda G(p;m),
\]
where $G$ is net public outlay and $\lambda\geq0$ its monetary social cost. Determine the welfare effect of adequacy standards, continuity rules, and risk adjustment after endogenous premiums and enrollment reequilibrate. Establish what quasi-random network changes identify access/health effects separately from selection into plans and renegotiation of provider prices. Report $\varepsilon$-optimal feasible regulations and bounds over models that fit plan-choice/claims data equally well. A narrow versus broad label is insufficient to define the counterfactual policy.

\subsection*{Interpretation and scope}
The relevant patient population includes people who leave a plan or become uninsured, not only continuously enrolled claimants. Insurers may exclude a high-quality provider to discourage costly patients, so lower average observed spending can reflect selection rather than efficiency. Conversely, selective contracting can improve bargaining and lower premiums; the welfare gain must be measured after access losses and capacity adjustments. Premiums and provider payments are transfers within the welfare boundary but have distributional and utilization effects; avoid counting them as resource savings. Nonrandom network changes can be caused by anticipated demand or provider quality changes, requiring explicit identification assumptions. The extension compares implementable network regulation in joint bargaining/selection equilibria with clinical outcomes and travel costs.

\subsection*{Source audit and status}
All three original sources are identifiable. Shepard (2016) is NBER WP22600; its published AER version is 2022, not 2016. [S1] provides network-model antecedents and [S3] a particular plan evaluation. These studies do not establish a uniquely optimal breadth rule for every market or certify current unresolved status.

\subsection*{References}
\begin{enumerate}[label={[S\arabic*]},leftmargin=*,itemsep=0.6em]
\item Katherine Ho (2009). Insurer-Provider Networks in the Medical Care Market. American Economic Review 99(1), 393--430. DOI: 10.1257/aer.99.1.393.
\url{https://www.aeaweb.org/articles?id=10.1257/aer.99.1.393}
\newline\emph{Locator and support limits:} Abstract and model: endogenous network contracting; network breadth is multidimensional.
\item Mark Shepard (2016). Hospital Network Competition and Adverse Selection: Evidence from the Massachusetts Health Insurance Exchange. NBER Working Paper 22600. DOI: 10.3386/w22600. Published version: American Economic Review 112(2), 578--615 (2022), DOI 10.1257/aer.20201453.
\url{https://www.nber.org/papers/w22600}
\newline\emph{Locator and support limits:} Abstract: networks and adverse selection. Preserve 2016 as working-paper year; the journal article appeared in 2022.
\item Jonathan Gruber and Robin McKnight (2016). Controlling Health Care Costs through Limited Network Insurance Plans: Evidence from Massachusetts State Employees. American Economic Journal: Economic Policy 8(2), 219--250. DOI: 10.1257/pol.20140335.
\url{https://www.aeaweb.org/articles?id=10.1257/pol.20140335}
\newline\emph{Locator and support limits:} Abstract and evaluation: limited-network plans for Massachusetts state employees; clinical/access transport needs further evidence.
\end{enumerate}

\subsection*{Common conventions: Source mathematical conventions}

Unless an entry states otherwise, agent and firm sets are finite. State and action spaces are measurable, strategies and outcome maps are measurable, probability laws are specified, and expectations exist for the targets under discussion. Infinite-horizon discount factors lie in $(0,1)$ and discounted objectives are finite. Norms, tolerances and complexity input sizes must be specified for computational comparisons. Constants or primitives that appear locally are inputs, not empirical facts asserted by this catalogue.

\textbf{Identification.} A statistical model $m\in\mathcal M$ implies an observable law $P_m$ and target $\theta(m)$. At law $P$, its identified set is
\[
\Theta(P;\mathcal M)=\{\theta(m):m\in\mathcal M,\ P_m=P\}.
\]
Point identification holds when this set is a singleton, or equivalently when $P_{m_1}=P_{m_2}$ implies $\theta(m_1)=\theta(m_2)$ for all compatible models. Sharp bounds mean bounds attained, or approached when the boundary is not attained, by compatible models. Writing this set is a definition, not a solution: the substantive task is to establish informative restrictions and derive or estimate the set. An unrestricted-confounding model may give only trivial bounds.

\textbf{Inference.} Identification, estimability and valid uncertainty quantification are separate questions. Fix $\alpha\in(0,1)$. A confidence procedure $C_n$ over a declared law class $\mathcal P$ has asymptotic uniform coverage at level $1-\alpha$ when
\[
\liminf_{n\to\infty}\inf_{P\in\mathcal P}\Pr_P\{\theta(P)\in C_n\}\geq1-\alpha.
\]
For set-identified targets the coverage event must be defined separately, for example coverage of the full identified set. If an entry uses identification-indexed models instead of uniquely indexed laws, the same coverage convention is applied over those models. Pointwise limits do not establish uniform validity, and asymptotic validity does not guarantee finite-sample precision. Sample sizes, dependence structures and weak-identification sequences are part of each application.

\textbf{Counterfactuals.} $Y_i(a)$ is a potential outcome under a specified intervention, including its timing, implementation and versions. It is not $Y_i$ conditional on voluntarily choosing $a$. A full intervention vector permits interference; shorthand scalar treatment notation does not silently impose no interference. Assignment, selection, attrition, anticipation and measurement must be specified. A claim about an instrument's local effect is not automatically a population policy effect.

\textbf{Economic equilibrium.} A counterfactual model includes best responses, constraints, feasibility, market clearing, state transitions and expectations consistent with the stated information. When several equilibria exist, either a declared selector is an input or the target is set-valued. Comparisons across policy regimes must use the stated selection convention. Reduced-form relationships are not presumed invariant after an equilibrium-changing intervention.

\textbf{Optimization and welfare.} Optimization is over the stated feasible set. If attainment is not established, $\sup$ or $\inf$ and $\varepsilon$-optimal choices replace an asserted maximizer or minimizer. For example, with value $V=\sup_{a\in\mathcal A}W(a)<\infty$, an $\varepsilon$-optimal set is $\{a\in\mathcal A:W(a)\geq V-\varepsilon\}$ for $\varepsilon>0$. Benefits, costs, risk and health or environmental outcomes can be aggregated only using declared units and conversion weights. Interpersonal utility weights, the treatment of fiscal transfers and foreign profits, discounting, and the behavioral welfare criterion are explicit inputs. A comparative-static derivative additionally requires existence and appropriate differentiability of the selected counterfactual.

\textbf{Sources.} Local labels [S1], [S2], and so forth restart in every question. Locators identify the relevant abstract, model, section or result at the level actually checked; a bibliographic or abstract-level verification is not represented as inspection of an entire proof. Working papers are distinguished from later publications and changing versions. Source summaries are paraphrased. All URL checks and editorial formulations refer to the audit date stated above.
% END ECONOMICS STATEMENT
\end{document}
